% Unidad U010. Espejo orientado y selección 3->1; redacción autónoma.

\chapter{Involution spéculaire et sélection locale \(3\to1\) de l'extension survivante}
\label{chap:espejo-seleccion-tres-uno}

La neuvième phase contient deux états locaux d'entrée. Une fois fixé
l'état orienté \(B_9\), la relation de frontière produit trois extensions
vers la première phase. Les trois partagent un axe d'autoéchelle et un agrégat
ternaire ; elles diffèrent par la répartition ordonnée de cet agrégat entre les branches de
clôture et de propagation. Un second horizon en conserve exactement une.

\section{Coordonnées symétrique et axiale}

Soit
\[
 \mathcal B:=(\mathbb F_3^6)^3.
\]
Un bloc visible est
\[
 B=(u^\pi,u^e,u^\varphi)\in\mathcal B.
\]
On définit
\[
 q(B):=u^\pi+u^e+u^\varphi,
 \qquad
 a(B):=u^\pi+u^e-u^\varphi.
\]
Toutes les opérations s'effectuent composante par composante dans
\(\mathbb F_3^6\).

\begin{theorem}[Reconstruction de l'axe et de l'agrégat]
\label{thm:espejo-reconstruccion-eje-agregado}
Les coordonnées \(q,a\) déterminent
\[
 \boxed{u^\varphi=2(q-a),}
 \qquad
 \boxed{u^\pi+u^e=q-u^\varphi.}
\]
\end{theorem}

\begin{proof}
Soustraire les définitions donne
\(q-a=2u^\varphi\). Dans \(\mathbb F_3\), \(2^{-1}=2\) ; par conséquent
\(u^\varphi=2(q-a)\). Substituer dans la définition de \(q\) récupère
l'agrégat restant.
\end{proof}

L'involution spéculaire est
\[
 \mathfrak c(u^\pi,u^e,u^\varphi)
 =(u^e,u^\pi,u^\varphi).
\]

\begin{proposition}[Invariants de l'involution]
\label{prop:espejo-invariantes}
\(\mathfrak c^2=\operatorname{id}\), y
\[
 q(\mathfrak cB)=q(B),
 \qquad
 a(\mathfrak cB)=a(B).
\]
En particulier, \(u^\varphi\) et \(u^\pi+u^e\) sont invariants.
\end{proposition}

\begin{proof}
Échanger deux fois les premières coordonnées récupère \(B\). Tant
\(q\) que \(a\) ne dépendent d'elles que par leur somme, de sorte que
l'échange les préserve. Le \cref{thm:espejo-reconstruccion-eje-agregado}
transfère cette invariance à l'axe et à l'agrégat.
\end{proof}

L'involution n'affirme pas que \(\pi=e\) ni que l'une soit le quotient de l'autre.
Elle affirme que les deux branches occupent des positions conjuguées par rapport à une
structure fixe.

\section{État orienté dans la phase \(9\) et données de frontière associées}

La projection visible de l'état sélectionné est
\[
 B_9=(100100\mid020112\mid010122).
\]
La relation induite sur les blocs visibles est notée
\[
 \mathcal M_9^B\subseteq\mathcal B\times\mathcal B.
\]
Les deux états locaux d'entrée de la table de phase précèdent cette
relation. Après sélection de l'état complet qui se projette sur \(B_9\), sa
fibre de sortie est de cardinal trois.

\begin{theorem}[Les trois extensions de la neuvième phase]
\label{thm:espejo-tres-extensiones}
L'image relationnelle de \(B_9\) contient exactement
\begin{align*}
B_{10}^{(0)}&=(101222\mid112221\mid112002),\\
B_{10}^{(1)}&=(102221\mid111222\mid112002),\\
B_{10}^{(2)}&=(111221\mid102222\mid112002).
\end{align*}
Les trois satisfont
\[
 q_{10}=022112,
 \qquad
 a_{10}=101111,
\]
et, par conséquent,
\[
 u_{10}^\varphi=112002,
 \qquad
 u_{10}^\pi+u_{10}^e=210110.
\]
\end{theorem}

\begin{proof}
La relation d'extension sur les états complets est énumérée dans la fibre de
\(B_9\) et produit les trois lignes indiquées. Additionner et soustraire leurs trois
composantes donne la même paire \((q_{10},a_{10})\) dans les trois cas. Le
\cref{thm:espejo-reconstruccion-eje-agregado} produit l'axe et l'agrégat
communs. Les deux premières composantes sont distinctes, si bien que les trois
lignes représentent des distributions différentes du même agrégat.
\end{proof}

Leurs publications entières en base trois sont
\[
 (279,352,365),
 \qquad
 (280,351,365),
 \qquad
 (282,350,365).
\]

\section{Compatibilité exacte des cylindres}

Un mot ternaire de longueur \(L\) et d'indice \(N\) détermine le cylindre
semi-ouvert
\[
 I_3(N,L)=\left[\frac{N}{3^L},\frac{N+1}{3^L}\right).
\]
Un préfixe de \(b\) blocs décimaux, d'entier associé \(D\), détermine
\[
 I_{1000}(D,b)=
 \left[\frac{D}{1000^b},\frac{D+1}{1000^b}\right).
\]

\begin{lemma}[Test entier d'intersection]
\label{lem:espejo-test-entero-cilindros}
Les cylindres \(I_3(N,L)\) et \(I_{1000}(D,b)\) se rencontrent si et seulement si
\[
 N1000^b<(D+1)3^L,
 \qquad
 (N+1)1000^b>D3^L.
\]
\end{lemma}

\begin{proof}
Deux intervalles semi-ouverts \([a,b)\) et \([c,d)\) s'intersectent si et seulement
si \(a<d\) et \(c<b\). Nous substituons les quatre extrémités et multiplions par
l'entier positif \(3^L1000^b\).
\end{proof}

Le test utilise uniquement l'arithmétique entière. Il ne dépend pas d'une tolérance en virgule
flottante.

\section{2e horizon et unicité}

Le bloc suivant est
\[
 B_{11}=(022212\mid110001\mid100021),
\]
et la frontière décimale correspondante est
\[
 (502,662,638).
\]

\begin{theorem}[Résolution locale \(3\to1\)]
\label{thm:espejo-seleccion-tres-uno}
Les trois extensions du \cref{thm:espejo-tres-extensiones} sont compatibles
à profondeur un. En exigeant simultanément \(B_{11}\) et la frontière
\((502,662,638)\), seule \(B_{10}^{(0)}\) se prolonge comme le même état
complet. Par conséquent,
\[
 \boxed{
 \#\operatorname{Surv}^{\rm cal}_1(B_9)=3,
 \qquad
 \#\operatorname{Surv}^{\rm cal}_2(B_9)=1.}
\]
\end{theorem}

\begin{proof}
Nous appliquons le \cref{lem:espejo-test-entero-cilindros} à chacun des trois
canaux. Au premier horizon, les six inégalités de chaque candidat sont
satisfaites. Après ajout de \(B_{11}\), les inégalités des trois canaux
se maintiennent pour \(B_{10}^{(0)}\). Dans \(B_{10}^{(1)}\) et
\(B_{10}^{(2)}\), la compatibilité du canal \(\varphi\) est maintenue, mais
les conditions des branches \(\pi\) et \(e\) échouent. L'énumération laisse
donc un unique état à profondeur deux.
\end{proof}

Cette preuve ne définit pas la connexion globale à partir d'un exemple local. C'est un
témoin explicite de la manière dont une relation multivaluée à un horizon peut
devenir univaluée en conservant la mémoire et en exigeant un prolongement plus grand.

\section{Retour de phase avec changement de bloc}

La première phase du tour initial est
\[
 B_1=(012222\mid121221\mid112202),
\]
tandis que la première phase du tour suivant est
\[
 B_{10}=(101222\mid112221\mid112002).
\]
Ainsi,
\[
 g(1)=g(10)=1,
 \qquad
 B_1\neq B_{10}.
\]
La phase revient, mais l'axe, l'agrégat, la répartition et la mémoire de
quotient ont évolué.

\section{Diagramme de l'involution spéculaire \(\pi/e\), de l'axe fixe \(\varphi\) et de la sélection locale \(3\to1\) par prolongeabilité}

\begin{figure}[htbp]
\centering
\resizebox{0.98\textwidth}{!}{%
\begin{tikzpicture}[
 node distance=10mm and 15mm,
 every node/.style={font=\small,align=center},
 state/.style={draw=black,rounded corners,inner sep=5pt},
 inv/.style={draw=black,double,rounded corners,inner sep=5pt},
 arr/.style={-{Latex[length=2mm]},semithick}
]
 \node[state] (pi) {$u^\pi$\\branche de clôture};
 \node[state,right=70mm of pi] (e) {$u^e$\\branche de propagation};
 \coordinate (mirroraxis) at ($(pi)!0.5!(e)$);
 \node[inv,above=9mm of mirroraxis] (phi) {$u^\varphi$\\eje fijo};
 \draw[arr,<->] (pi) -- node[below] {$\mathfrak c$} (e);
 \node[above=8mm of phi] (aggregate)
   {$u^\pi+u^e$ conservé dans la fibre};
 \draw[dashed] (phi.north) -- (aggregate.south);
 \node[state,below=20mm of mirroraxis] (b9) {phase $9$ : $B_9$\\deux entrées locales};
 \node[state,below left=18mm and 25mm of b9] (b100) {$B_{10}^{(0)}$};
 \node[state,below=18mm of b9] (b101) {$B_{10}^{(1)}$};
 \node[state,below right=18mm and 25mm of b9] (b102) {$B_{10}^{(2)}$};
 \node[state,below=18mm of b101] (b11) {$B_{11}$\\unique prolongement};
 \draw[arr] (b9) -- (b100);
 \draw[arr] (b9) -- (b101);
 \draw[arr] (b9) -- (b102);
 \draw[arr] (b100) -- (b11);
 \draw[arr,densely dotted] (b101) -- node[right] {obstruée} (b11);
\draw[arr,densely dotted] (b102) -- node[right] {obstruée} (b11);
\end{tikzpicture}
}
\caption{Le miroir conserve l'axe \(u^\varphi\) et l'agrégat
\(u^\pi+u^e\). Le retour de la neuvième phase ouvre trois extensions visibles ;
la mémoire du second horizon en conserve une. Le tour ne réinitialise pas
l'état.}
\label{fig:espejo-tres-uno-arrastre}
\end{figure}

\section{Mémoire résidu--quotient pendant le miroir}

La récurrence
\[
 x_kA=u_k+3x_{k+1}
\]
de la \cref{eq:bockstein-hensel-local} s'applique à chacun des trois
canaux. La partie \(u_k\) publie le bloc visible ; \(x_{k+1}\) conserve
le quotient. Comme \(A\in\operatorname{GL}_6(\mathbb Z)\), la paire récupère
l'état précédent. Le miroir peut échanger les branches visibles sans
éliminer leurs quotients de provenance.

\section{Obstructions de l'involution spéculaire \(\pi/e\) et de la sélection locale \(3\to1\) : invariants, horizon de profondeur \(2\) et mémoire de Hensel}

\begin{criterion}[Critères de réfutation du miroir]
La construction est réfutée si :
\begin{enumerate}
 \item l'involution modifie \(u^\varphi\) ou \(u^\pi+u^e\) ;
 \item les deux états locaux d'entrée sont confondus avec les trois sorties ;
 \item l'une des trois sorties possède un \((q_{10},a_{10})\) différent ;
 \item le test des cylindres utilise une tolérance flottante ;
 \item plus d'une sortie survit au second horizon déclaré ;
 \item le retour de phase est interprété comme l'égalité \(B_1=B_{10}\) ;
 \item l'échange visible efface les quotients de Hensel.
\end{enumerate}
\end{criterion}

Le miroir et la sélection locale sont maintenant construits. Les unités
suivantes développeront les lacunes sturmiennes puis les trois
biographies analytiques complètes.
